% =====================================================================
%  Analítica de Datos — Sesión 05
%  Universidad Cooperativa de Colombia — Ingeniería de Software
%  Autor: Luis Miguel Gómez Meneses
%  Estructura: CONTINUIDAD -> TEORÍA -> LABORATORIO (4 h)
%  Compilar con pdfLaTeX en Overleaf
% =====================================================================

\documentclass[aspectratio=169,11pt]{beamer}

\usepackage[T1]{fontenc}
\usepackage[utf8]{inputenc}
\usepackage[spanish,es-noshorthands,es-tabla]{babel}
\usepackage{lmodern}
\usepackage{xcolor}
\usepackage{array}
\usepackage{booktabs}
\usepackage{listings}
\usepackage{upquote}
\usepackage{tikz}
\usetikzlibrary{arrows.meta,positioning,calc,shapes.geometric}

\usetheme{default}
\usefonttheme{professionalfonts}
\usepackage{graphicx}
\setbeamertemplate{navigation symbols}{}

% ------------------------------ Paleta -------------------------------

\definecolor{cPrim}{HTML}{14484F}
\definecolor{cAcc}{HTML}{C2571A}
\definecolor{cInk}{HTML}{1C2B2D}
\definecolor{cSoft}{HTML}{EDF2F2}
\definecolor{cGray}{HTML}{6B7A7C}
\definecolor{cCode}{HTML}{F6F4EF}
\definecolor{cOut}{HTML}{F2F2F2}

\setbeamercolor{normal text}{fg=cInk}
\setbeamercolor{structure}{fg=cPrim}
\setbeamercolor{frametitle}{fg=cPrim,bg=white}
\setbeamercolor{itemize item}{fg=cAcc}
\setbeamercolor{itemize subitem}{fg=cGray}
\setbeamercolor{block title}{fg=white,bg=cPrim}
\setbeamercolor{block body}{fg=cInk,bg=cSoft}
\setbeamercolor{block title alerted}{fg=white,bg=cAcc}
\setbeamercolor{block body alerted}{fg=cInk,bg=cAcc!10}
\setbeamercolor{block title example}{fg=white,bg=cGray}
\setbeamercolor{block body example}{fg=cInk,bg=cGray!12}

\setbeamerfont{frametitle}{size=\large,series=\bfseries}
\setbeamerfont{block title}{size=\small,series=\bfseries}
\setbeamerfont{block body}{size=\small}

\setbeamertemplate{itemize item}
  {\raisebox{0.15ex}{\rule{1.1ex}{1.1ex}}}

\setbeamertemplate{itemize subitem}{\textendash}
\setbeamertemplate{blocks}[rounded][shadow=false]

% ------------------- Modo laboratorio (marca visual) -----------------

\newcommand{\ruleclr}{cAcc}
\newcommand{\marcaframe}{}

\newcommand{\modolab}[1]{%
  \setbeamercolor{frametitle}{fg=cAcc}%
  \setbeamercolor{background canvas}{bg=cAcc!8}%
  \renewcommand{\ruleclr}{cPrim}%
  \renewcommand{\marcaframe}{#1}%
}

\newcommand{\modoteoria}{%
  \setbeamercolor{frametitle}{fg=cPrim}%
  \setbeamercolor{background canvas}{bg=white}%
  \renewcommand{\ruleclr}{cAcc}%
  \renewcommand{\marcaframe}{}%
}

% --------------------------- Título de marco --------------------------

\setbeamertemplate{frametitle}{%
  \vspace*{0.55em}%
  \begin{beamercolorbox}
    [wd=\paperwidth,leftskip=0.9cm,rightskip=0.9cm]{frametitle}%
    \usebeamerfont{frametitle}\insertframetitle\par
    \vspace{0.25em}%
    {\color{\ruleclr}\rule{2.1cm}{1.2pt}}\hfill
    {\scriptsize\bfseries\color{cGray}\marcaframe}\par
  \end{beamercolorbox}%
  \vspace{0.2em}%
}

% ------------------------------ Pie ----------------------------------

\setbeamertemplate{footline}{%
  \hbox{%
  \begin{beamercolorbox}
    [wd=\paperwidth,ht=2.4ex,dp=1.4ex,
     leftskip=0.9cm,rightskip=0.9cm]{}%
    {\tiny\color{cGray}\insertsection\hfill\insertframenumber}%
  \end{beamercolorbox}}%
}

% ------------------------------ TikZ ---------------------------------

\tikzset{
  nodo/.style={
    rounded corners=2pt,
    draw=cPrim!55,
    fill=cSoft,
    align=center,
    inner sep=3.5pt,
    font=\scriptsize,
    text=cInk,
    minimum height=0.85cm
  },
  nodoA/.style={
    nodo,
    draw=cAcc!70,
    fill=cAcc!12
  },
  nodoD/.style={
    nodo,
    draw=cGray!60,
    fill=white
  },
  flecha/.style={
    -{Stealth[length=1.9mm]},
    draw=cPrim!75,
    line width=0.7pt
  },
  flechaA/.style={
    -{Stealth[length=1.9mm]},
    draw=cAcc!85,
    line width=0.7pt
  },
  etiq/.style={
    font=\scriptsize\itshape,
    text=cGray
  }
}

% ---------------------------- Listings -------------------------------

\lstdefinestyle{py}{
  language=Python,
  basicstyle=\ttfamily\small,
  keywordstyle=\color{cPrim}\bfseries,
  stringstyle=\color{cAcc},
  commentstyle=\color{cGray}\itshape,
  showstringspaces=false,
  columns=fullflexible,
  keepspaces=true,
  breaklines=false,
  numbers=none,
  frame=leftline,
  framerule=2pt,
  rulecolor=\color{cPrim!45},
  backgroundcolor=\color{cCode},
  xleftmargin=8pt,
  framexleftmargin=8pt,
  framexrightmargin=4pt,
  aboveskip=5pt,
  belowskip=5pt,
  literate=%
    {á}{{\'a}}1
    {é}{{\'e}}1
    {í}{{\'i}}1
    {ó}{{\'o}}1
    {ú}{{\'u}}1
    {Á}{{\'A}}1
    {É}{{\'E}}1
    {Í}{{\'I}}1
    {Ó}{{\'O}}1
    {Ú}{{\'U}}1
    {ñ}{{\~n}}1
    {Ñ}{{\~N}}1
    {¿}{{?`}}1
    {¡}{{!`}}1
}

\lstdefinestyle{salida}{
  basicstyle=\ttfamily\scriptsize\color{cInk},
  showstringspaces=false,
  columns=fullflexible,
  keepspaces=true,
  numbers=none,
  frame=leftline,
  framerule=2pt,
  rulecolor=\color{cGray!45},
  backgroundcolor=\color{cOut},
  xleftmargin=8pt,
  framexleftmargin=8pt,
  aboveskip=5pt,
  belowskip=5pt,
  literate=%
    {á}{{\'a}}1
    {é}{{\'e}}1
    {í}{{\'i}}1
    {ó}{{\'o}}1
    {ú}{{\'u}}1
    {ñ}{{\~n}}1
}

\lstset{style=py}

% ---------------------------- Utilidades -----------------------------

\newcommand{\kicker}[1]{%
  {\scriptsize\bfseries\color{cAcc}\MakeUppercase{#1}}%
  \par\vspace{0.2em}%
}

\newcommand{\destaca}[1]{%
  {\color{cPrim}\bfseries #1}%
}

\newcommand{\pregunta}[1]{%
  {\color{cAcc}\bfseries #1}%
}

\newenvironment{cerrado}[1]
  {\begin{block}{#1}
   \begin{minipage}{\linewidth}
   \raggedright}
  {\end{minipage}
   \end{block}}

% ---------------------------- Metadatos -------------------------------

\title{Analítica de Datos}
\author{Luis Miguel Gómez Meneses}
\institute{Universidad Cooperativa de Colombia}
\date{}

% =====================================================================
%  DOCUMENTO
% =====================================================================

\begin{document}

\setlength{\parskip}{0pt}

% =====================================================================
\section{Apertura}
% =====================================================================

% ---------------------------------------------------------------------
% Diapositiva 1 — Portada
% ---------------------------------------------------------------------

\begin{frame}[plain]

  \vspace{1.15cm}

  \hspace*{0.7cm}%
  \begin{minipage}{0.90\textwidth}

    \raggedright

    {\color{cAcc}\rule{2.6cm}{2.2pt}}\par

    \vspace{0.6em}

    {\Large\color{cGray}
      Analítica de Datos
    }\par

    \vspace{0.35em}

    {\LARGE\bfseries\color{cPrim}
      Sesión 05. Regresión lineal simple y múltiple
    }\par

    \vspace{0.9em}

    {\small\color{cInk}
      Modelos con una o varias variables predictoras,\\
      interpretación de coeficientes y evaluación
    }\par

    \vspace{1.5em}

    {\small\bfseries
      Luis Miguel Gómez Meneses
    }\par

    \vspace{0.3em}

    {\footnotesize\color{cGray}
      Universidad Cooperativa de Colombia\\
      Ingeniería de Software --- Séptimo semestre
    }\par

  \end{minipage}

\end{frame}

% ---------------------------------------------------------------------
% Diapositiva 2 — Cómo trabajaremos hoy
% ---------------------------------------------------------------------

\begin{frame}{Cómo trabajaremos hoy}

  \vspace{0.1em}

  \begin{center}

  \small
  \renewcommand{\arraystretch}{1.18}

  \begin{tabular}{@{}>{\raggedright\arraybackslash}p{1.9cm}
                    >{\raggedright\arraybackslash}p{7.2cm}
                    >{\raggedleft\arraybackslash}p{1.4cm}@{}}

    \toprule

    \textbf{Bloque} &
    \textbf{Contenido} &
    \textbf{Tiempo}\\

    \midrule

    Evaluación &
    Sustentaciones y cierre del primer 20\,\% &
    60 min\\

    Teoría 1 &
    Regresión lineal simple: ecuación, predicción y residuos &
    30 min\\

    \textcolor{cAcc}{\textbf{Lab 13}} &
    \textcolor{cAcc}{Ajuste e interpretación de un modelo simple} &
    \textcolor{cAcc}{35 min}\\

    Teoría 2 &
    Regresión lineal múltiple: predictores y coeficientes &
    30 min\\

    \textcolor{cAcc}{\textbf{Lab 14}} &
    \textcolor{cAcc}{Construcción de un modelo múltiple con Python} &
    \textcolor{cAcc}{45 min}\\

    Evaluación &
    Comparación de modelos: $R^2$, MAE, MSE y RMSE &
    25 min\\

    Cierre &
    Interpretación, conclusiones y actividad final &
    15 min\\

    \bottomrule

  \end{tabular}

  \end{center}

  \vspace{0.35em}

  \begin{center}
    \small\color{cGray}
    De una variable predictora a modelos con múltiples predictores.
  \end{center}

\end{frame}

% =====================================================================
\section{De la inferencia a la predicción}
\modoteoria
% =====================================================================

% ---------------------------------------------------------------------
% Diapositiva 3 — De la sesión anterior a esta sesión
% ---------------------------------------------------------------------

\begin{frame}{De la sesión anterior a esta sesión}

  \begin{center}
    \begin{minipage}{0.90\textwidth}

      \begin{cerrado}{En la sesión anterior}

        Utilizamos una muestra para estimar un parámetro desconocido
        de la población:

        \begin{itemize}
          \item calculamos la media muestral;
          \item determinamos el error estándar;
          \item construimos un intervalo de confianza;
          \item analizamos el efecto del tamaño de la muestra.
        \end{itemize}

      \end{cerrado}

      \vspace{0.35em}

      \begin{cerrado}{En esta sesión}

        Utilizaremos una o varias variables para explicar y predecir
        el comportamiento de una variable respuesta:

        \begin{itemize}
          \item con un predictor: \destaca{regresión lineal simple};
          \item con varios predictores: \destaca{regresión lineal múltiple}.
        \end{itemize}

      \end{cerrado}

      \vspace{0.3em}

      \begin{center}
        \pregunta{¿Cómo construimos un modelo que represente la relación
        entre las variables?}
      \end{center}

    \end{minipage}
  \end{center}

\end{frame}

% =====================================================================
\section{Regresión lineal}
\modoteoria
% =====================================================================

% ---------------------------------------------------------------------
% Diapositiva 4 — ¿Qué es la regresión lineal simple?
% ---------------------------------------------------------------------

\begin{frame}{¿Qué es la regresión lineal simple?}

  \begin{center}
    \begin{minipage}{0.90\textwidth}

      \kicker{DE LA RELACIÓN A LA PREDICCIÓN}

      \vspace{0.4em}

      La \textbf{regresión lineal simple} es un método estadístico
      que permite representar la relación lineal entre una variable
      predictora y una variable respuesta mediante una ecuación.

      \vspace{0.7em}

      En este modelo intervienen:

      \begin{itemize}
        \item una variable predictora, utilizada para explicar o predecir;
        \item una variable respuesta, cuyo valor queremos estimar.
      \end{itemize}

      \vspace{0.6em}

      \begin{cerrado}{Idea clave}

        Se denomina \textbf{simple} porque utiliza una sola variable
        predictora para estimar el comportamiento de la variable respuesta.

      \end{cerrado}

      \vspace{0.5em}

      \begin{center}
        \textbf{Un predictor}
        $\quad\longrightarrow\quad$
        \textbf{Modelo lineal}
        $\quad\longrightarrow\quad$
        \textbf{Valor estimado}
      \end{center}

    \end{minipage}
  \end{center}

\end{frame}

% ---------------------------------------------------------------------
% Diapositiva 5 — Variable predictora y variable respuesta
% ---------------------------------------------------------------------

\begin{frame}{Variable predictora y variable respuesta}

  \begin{center}
    \begin{minipage}{0.90\textwidth}

      \kicker{VARIABLES CON FUNCIONES DIFERENTES}

      \vspace{0.4em}

      En un modelo de regresión intervienen una variable respuesta
      y una o varias variables predictoras.

      \vspace{0.6em}

      \begin{columns}[T,onlytextwidth]

        \begin{column}{0.47\textwidth}

          \begin{cerrado}{Variable predictora \(X\)}

            Es una variable conocida que utilizamos para explicar
            o estimar el comportamiento de la variable respuesta.

            \vspace{0.4em}

            También se denomina:

            \begin{itemize}
              \item variable independiente;
              \item variable explicativa.
            \end{itemize}

          \end{cerrado}

        \end{column}

        \begin{column}{0.47\textwidth}

          \begin{cerrado}{Variable respuesta \(Y\)}

            Es la variable que queremos explicar o estimar mediante
            el modelo.

            \vspace{0.4em}

            También se denomina:

            \begin{itemize}
              \item variable dependiente;
              \item variable objetivo.
            \end{itemize}

          \end{cerrado}

        \end{column}

      \end{columns}

      \vspace{0.45em}

      \begin{center}
        \destaca{Regresión simple: una \(X\)}
        \qquad
        \destaca{Regresión múltiple: varias \(X\)}
      \end{center}

      \vspace{0.2em}

      \begin{center}
        \pregunta{En ambos casos utilizamos los predictores para estimar \(Y\).}
      \end{center}

    \end{minipage}
  \end{center}

\end{frame}

% ---------------------------------------------------------------------
% Diapositiva 6 — Un problema de regresión
% ---------------------------------------------------------------------

\begin{frame}{Un problema de regresión}

  \begin{center}
    \begin{minipage}{0.90\textwidth}

      \kicker{DE UNA PREGUNTA A DOS VARIABLES}

      \vspace{0.4em}

      Una empresa quiere analizar si la inversión mensual en publicidad
      puede utilizarse para estimar sus ventas.

      \vspace{0.7em}

      \begin{columns}[T,onlytextwidth]

        \begin{column}{0.46\textwidth}

          \begin{cerrado}{Variable predictora}

            \[
              X=\text{inversión en publicidad}
            \]

            Es la información conocida que utilizaremos para realizar
            la estimación.

          \end{cerrado}

        \end{column}

        \begin{column}{0.46\textwidth}

          \begin{cerrado}{Variable respuesta}

            \[
              Y=\text{ventas}
            \]

            Es la variable cuyo valor queremos explicar o predecir.

          \end{cerrado}

        \end{column}

      \end{columns}

      \vspace{0.7em}

      \begin{cerrado}{Pregunta de análisis}

        \begin{center}
          \pregunta{¿Cuánto esperamos vender para una determinada
          inversión en publicidad?}
        \end{center}

      \end{cerrado}

    \end{minipage}
  \end{center}

\end{frame}

% ---------------------------------------------------------------------
% Diapositiva 7 — El diagrama de dispersión
% ---------------------------------------------------------------------

\begin{frame}{El diagrama de dispersión}

  \begin{center}
    \begin{minipage}{0.92\textwidth}

      \kicker{PRIMERO OBSERVAMOS LA RELACIÓN}

      Antes de construir el modelo, representamos cada observación
      mediante un punto con coordenadas \((X,Y)\).

      \vspace{0.0em}

      \begin{columns}[T,onlytextwidth]

        % -------------------- Gráfica --------------------
        \begin{column}{0.52\textwidth}

          \begin{center}

          \begin{tikzpicture}[x=0.62cm,y=0.42cm]

            % Cuadrícula
            \draw[step=1cm,cGray!20,very thin]
              (0,0) grid (6,5);

            % Eje horizontal
            \draw[-{Stealth[length=2mm]},line width=0.8pt,cPrim]
              (0,0) -- (6.4,0)
              node[right,font=\scriptsize] {Publicidad \(X\)};

            % Eje vertical
            \draw[-{Stealth[length=2mm]},line width=0.8pt,cPrim]
              (0,0) -- (0,5.5)
              node[above,font=\scriptsize] {Ventas \(Y\)};

            % Observaciones
            \fill[cAcc] (0.6,0.5) circle (2pt);
            \fill[cAcc] (1.0,0.9) circle (2pt);
            \fill[cAcc] (1.6,1.2) circle (2pt);
            \fill[cAcc] (2.0,1.6) circle (2pt);
            \fill[cAcc] (2.5,1.4) circle (2pt);
            \fill[cAcc] (3.1,2.5) circle (2pt);
            \fill[cAcc] (3.6,2.8) circle (2pt);
            \fill[cAcc] (4.2,3.4) circle (2pt);
            \fill[cAcc] (4.7,4.0) circle (2pt);
            \fill[cAcc] (5.3,4.5) circle (2pt);
            \fill[cAcc] (5.7,4.9) circle (2pt);

          \end{tikzpicture}

          \end{center}

        \end{column}

        % ---------------- Interpretación ----------------
        \begin{column}{0.44\textwidth}

          \begin{cerrado}{¿Qué observamos?}

            \begin{itemize}
              \setlength{\itemsep}{0.2em}

              \item Los puntos presentan una tendencia ascendente.

              \item Cuando \(X\) aumenta, \(Y\) tiende a aumentar.

              \item La relación parece aproximadamente lineal.

              \item Los puntos no forman una recta perfecta.

            \end{itemize}

          \end{cerrado}

        \end{column}

      \end{columns}

      \vspace{0.15em}

      \begin{center}
        \pregunta{¿Qué recta representa mejor la tendencia de los datos?}
      \end{center}

    \end{minipage}
  \end{center}

\end{frame}
% ---------------------------------------------------------------------
% Diapositiva 8 — La ecuación de regresión lineal simple
% ---------------------------------------------------------------------

\begin{frame}{La ecuación de regresión lineal simple}

  \begin{center}
    \begin{minipage}{0.88\textwidth}

      \kicker{UNA RECTA PARA REALIZAR ESTIMACIONES}

      La relación entre la variable predictora \(X\) y la variable
      respuesta \(Y\) se representa mediante:

      \vspace{0.0em}

      \[
        \hat{Y}=\beta_0+\beta_1X
      \]

      \vspace{0.0em}

      \begin{cerrado}{Componentes del modelo}

        \small

        \begin{itemize}
          \setlength{\itemsep}{0.15em}

          \item \(X\): valor conocido de la variable predictora.

          \item \(\hat{Y}\): valor de \(Y\) estimado por el modelo.

          \item \(\beta_0\): intercepto o punto de partida de la recta.

          \item \(\beta_1\): pendiente de la recta. Indica cuánto cambia
          \(\hat{Y}\) cuando \(X\) aumenta una unidad.

        \end{itemize}

      \end{cerrado}

      \vspace{0.0em}

      \begin{center}
        \destaca{El modelo utiliza un valor conocido de \(X\)
        para obtener una estimación \(\hat{Y}\).}
      \end{center}

    \end{minipage}
  \end{center}

\end{frame}
% ---------------------------------------------------------------------
% Diapositiva 9 — Valor observado y valor estimado
% ---------------------------------------------------------------------

\begin{frame}{Valor observado y valor estimado}

  \begin{center}
    \begin{minipage}{0.88\textwidth}

      \kicker{EL MODELO PRODUCE UNA ESTIMACIÓN}

      Para cada valor de \(X\), distinguimos entre el resultado registrado
      en los datos y el resultado calculado por el modelo.

      \vspace{0.5em}

      \begin{columns}[T,onlytextwidth]

        \begin{column}{0.46\textwidth}

          \begin{cerrado}{Valor observado \(Y\)}

            Es el valor real registrado en los datos.

            \vspace{0.5em}

            Si la empresa vendió realmente 27 unidades:

            \[
              Y=27
            \]

            Este valor proviene directamente de la observación.

          \end{cerrado}

        \end{column}

        \begin{column}{0.46\textwidth}

          \begin{cerrado}{Valor estimado \(\hat{Y}\)}

            Es el valor calculado mediante la ecuación del modelo.

            \vspace{0.5em}

            Si el modelo estimó 25 unidades:

            \[
              \hat{Y}=25
            \]

            Este valor constituye una predicción.

          \end{cerrado}

        \end{column}

      \end{columns}

      \vspace{0.0em}

      \begin{center}
        \destaca{\(Y\) proviene de los datos}
        \qquad
        \destaca{\(\hat{Y}\) proviene del modelo}
      \end{center}

    \end{minipage}
  \end{center}

\end{frame}
% ---------------------------------------------------------------------
% Diapositiva 10 — El residuo
% ---------------------------------------------------------------------

\begin{frame}{El residuo}

  \begin{center}
    \begin{minipage}{0.86\textwidth}

      \kicker{DIFERENCIA ENTRE EL VALOR REAL Y LA ESTIMACIÓN}

      El residuo mide cuánto se aleja la predicción del valor
      realmente observado.

      \vspace{0.4em}

      \[
        e=Y-\hat{Y}
      \]

      \vspace{0.25em}

      \begin{cerrado}{Ejemplo}

        Si el valor observado es \(Y=27\) y el modelo estima
        \(\hat{Y}=25\), entonces:

        \[
          e=27-25=2
        \]

        El modelo estimó 2 unidades menos que el valor real.

      \end{cerrado}

      \vspace{0.45em}

      \begin{center}
        \destaca{El residuo compara el resultado real con la predicción
        del modelo.}
      \end{center}

    \end{minipage}
  \end{center}

\end{frame}
% ---------------------------------------------------------------------
% Diapositiva 11 — Interpretación del residuo
% ---------------------------------------------------------------------

\begin{frame}{Interpretación del residuo}

  \begin{center}
    \begin{minipage}{0.86\textwidth}

      \kicker{EL SIGNO INDICA LA DIRECCIÓN DEL ERROR}

      El signo del residuo permite saber si el modelo estimó un valor
      menor o mayor que el observado.

      \vspace{0.6em}

      \renewcommand{\arraystretch}{1.65}

      \begin{tabular}{@{}p{1.5cm}p{8.2cm}@{}}

        \toprule

        \textbf{Residuo} &
        \textbf{Interpretación}\\

        \midrule

        \(\boldsymbol{e>0}\) &
        El modelo estimó un valor menor que el observado.\\

        \(\boldsymbol{e<0}\) &
        El modelo estimó un valor mayor que el observado.\\

        \(\boldsymbol{e=0}\) &
        La estimación coincide con el valor observado.\\

        \bottomrule

      \end{tabular}

      \vspace{0.7em}

      \begin{cerrado}{Idea clave}

        Cuanto más cercano a cero sea el residuo, más próxima será
        la estimación al valor observado.

      \end{cerrado}

    \end{minipage}
  \end{center}

\end{frame}
% ---------------------------------------------------------------------
% Diapositiva 12 — La recta de mejor ajuste
% ---------------------------------------------------------------------

\begin{frame}{La recta de mejor ajuste}

  \begin{center}
    \begin{minipage}{0.84\textwidth}

      \kicker{UNA RECTA QUE REDUCE LOS ERRORES}

      Es posible dibujar diferentes rectas sobre los mismos datos,
      pero cada una produce estimaciones distintas.

      \vspace{0.3em}

      \begin{cerrado}{Criterio de ajuste}

        Para cada observación, el modelo calcula un residuo:

        \[
          e_i=Y_i-\hat{Y}_i
        \]

        La regresión lineal selecciona los coeficientes
        \(\beta_0\) y \(\beta_1\) que minimizan:

        \[
          \sum_{i=1}^{n}\left(Y_i-\hat{Y}_i\right)^2
        \]

        \vspace{0.15em}

        \destaca{La mejor recta es la que reduce, en conjunto,
        las diferencias entre los valores observados y estimados.}

      \end{cerrado}

    \end{minipage}
  \end{center}

\end{frame}
% ---------------------------------------------------------------------
% Diapositiva 13 — Coeficiente de determinación
% ---------------------------------------------------------------------

\begin{frame}{El coeficiente de determinación \(R^2\)}

  \begin{center}
    \begin{minipage}{0.84\textwidth}

      \kicker{CUÁNTA VARIABILIDAD EXPLICA EL MODELO}

      El coeficiente de determinación mide qué proporción de la
      variabilidad de \(Y\) explica el modelo de regresión.

      \[
        R^2
        =
        1-
        \frac{\sum_{i=1}^{n}(Y_i-\hat{Y}_i)^2}
        {\sum_{i=1}^{n}(Y_i-\bar{Y})^2}
      \]

      \begin{cerrado}{Interpretación}

        \begin{itemize}
          \setlength{\itemsep}{0.25em}

          \item Un valor cercano a \(1\) indica mayor capacidad explicativa.

          \item Un valor cercano a \(0\) indica que el modelo explica
          poca variabilidad de la respuesta.

          \item Por ejemplo, \(R^2=0.65\) significa que el modelo explica
          aproximadamente el \(65\,\%\) de la variabilidad observada.

        \end{itemize}

      \end{cerrado}

    \end{minipage}
  \end{center}

\end{frame}
% ---------------------------------------------------------------------
% Diapositiva 14 — Error absoluto medio
% ---------------------------------------------------------------------

\begin{frame}{El error absoluto medio}

  \begin{center}
    \begin{minipage}{0.84\textwidth}

      \kicker{EL TAMAÑO PROMEDIO DE LOS ERRORES}

      El error absoluto medio, o MAE, calcula el promedio de las
      diferencias absolutas entre los valores observados y estimados.

      \[
        MAE
        =
        \frac{1}{n}
        \sum_{i=1}^{n}
        |Y_i-\hat{Y}_i|
      \]

      \begin{cerrado}{Interpretación}

        \begin{itemize}
          \setlength{\itemsep}{0.25em}

          \item Ignora el signo de los residuos.

          \item Se expresa en las mismas unidades de la variable respuesta.

          \item Un MAE menor indica predicciones más cercanas a los
          valores observados.

        \end{itemize}

        Si \(MAE=2\), las predicciones se alejan de los valores reales
        aproximadamente 2 unidades, en promedio.

      \end{cerrado}

    \end{minipage}
  \end{center}

\end{frame}
% ---------------------------------------------------------------------
% Diapositiva 15 — MSE y RMSE
% ---------------------------------------------------------------------

\begin{frame}{Errores cuadráticos: MSE y RMSE}

  \begin{center}
    \begin{minipage}{0.84\textwidth}

      \kicker{PENALIZAR LOS ERRORES GRANDES}

      El error cuadrático medio eleva cada residuo al cuadrado:

      \[
        MSE
        =
        \frac{1}{n}
        \sum_{i=1}^{n}
        (Y_i-\hat{Y}_i)^2
      \]

      La raíz del error cuadrático medio se obtiene mediante:

      \[
        RMSE=\sqrt{MSE}
      \]

      \begin{cerrado}{Diferencia principal}

        \begin{itemize}
          \setlength{\itemsep}{0.2em}

          \item El MSE penaliza con mayor intensidad los errores grandes.

          \item El MSE queda expresado en unidades al cuadrado.

          \item El RMSE vuelve a las unidades originales de \(Y\), por
          lo que resulta más fácil de interpretar.

        \end{itemize}

      \end{cerrado}

    \end{minipage}
  \end{center}

\end{frame}
% =====================================================================
\section{Laboratorio 13}
\modolab{LABORATORIO 13}
% =====================================================================

% ---------------------------------------------------------------------
% Diapositiva 13 — Laboratorio de regresión lineal simple
% ---------------------------------------------------------------------

\begin{frame}{Regresión lineal simple con Python}

  \begin{center}
    \begin{minipage}{0.86\textwidth}

      \kicker{DE LOS DATOS AL MODELO}

      En este laboratorio construiremos un modelo para estimar las
      ventas de una empresa a partir de su inversión en publicidad.

      \vspace{0.4em}

      \begin{cerrado}{Pregunta de análisis}

        \begin{center}
          \pregunta{¿Podemos utilizar la inversión en publicidad
          para estimar las ventas?}
        \end{center}

      \end{cerrado}

      \vspace{0.35em}

      \begin{cerrado}{Ruta de trabajo}

        \begin{enumerate}
          \setlength{\itemsep}{0.2em}

          \item Explorar las variables del conjunto de datos.
          \item Construir el diagrama de dispersión.
          \item Ajustar el modelo de regresión lineal.
          \item Interpretar el intercepto y la pendiente.
          \item Obtener predicciones y calcular residuos.

        \end{enumerate}

      \end{cerrado}

    \end{minipage}
  \end{center}

\end{frame}


% =====================================================================
% FIN DEL LABORATORIO 13
% REGRESO AL MODO TEORÍA
% =====================================================================

\section{Regresión lineal múltiple}
\modoteoria


% ---------------------------------------------------------------------
% Diapositiva 18 — ¿Qué es la regresión lineal múltiple?
% ---------------------------------------------------------------------

\begin{frame}{¿Qué es la regresión lineal múltiple?}

  \begin{center}
    \begin{minipage}{0.90\textwidth}

      \small

      \kicker{VARIAS VARIABLES PARA ESTIMAR UNA RESPUESTA}

      La \textbf{regresión lineal múltiple} permite estimar una
      variable respuesta utilizando dos o más variables predictoras.

      \vspace{0.3em}

      Su ecuación general es:

      \vspace{-0.3em}

      \[
        \hat{Y}
        =
        \beta_0
        + \beta_1X_1
        + \beta_2X_2
        + \cdots
        + \beta_pX_p
      \]

      \vspace{-0.4em}

      \begin{cerrado}{Componentes del modelo}

        \begin{itemize}
          \setlength{\itemsep}{0.2em}

          \item \(\hat{Y}\): valor estimado de la variable respuesta.

          \item \(\beta_0\): intercepto o valor inicial del modelo.

          \item \(X_1,X_2,\ldots,X_p\): variables predictoras.

          \item \(\beta_1,\beta_2,\ldots,\beta_p\): cambios estimados
                en \(Y\) asociados con cada variable predictora.

        \end{itemize}

      \end{cerrado}

      \vspace{0.3em}

      \begin{center}
        \destaca{Cada coeficiente se interpreta manteniendo constantes
        las demás variables del modelo.}
      \end{center}

    \end{minipage}
  \end{center}

\end{frame}


% ---------------------------------------------------------------------
% Diapositiva 19 — Modelo de regresión múltiple para las ventas
% ---------------------------------------------------------------------

\begin{frame}{Modelo de regresión múltiple para las ventas}

  \begin{center}
    \begin{minipage}{0.90\textwidth}

      \small

      \kicker{TRES VARIABLES PREDICTORAS}

      Utilizaremos las inversiones en televisión, radio y prensa
      para estimar las ventas:

      \vspace{-0.2em}

      \[
        \widehat{\text{sales}}
        =
        \beta_0
        + \beta_1(\text{TV})
        + \beta_2(\text{radio})
        + \beta_3(\text{newspaper})
      \]

      \vspace{-0.4em}

      \begin{columns}[T,onlytextwidth]

        \begin{column}{0.47\textwidth}

          \begin{cerrado}{Variables predictoras}

            \begin{itemize}
              \setlength{\itemsep}{0.2em}

              \item \(X_1=\text{TV}\)
              \item \(X_2=\text{radio}\)
              \item \(X_3=\text{newspaper}\)

            \end{itemize}

          \end{cerrado}

        \end{column}

        \begin{column}{0.47\textwidth}

          \begin{cerrado}{Variable respuesta}

            \[
              Y=\text{sales}
            \]

            Ventas del producto expresadas en miles de unidades.

          \end{cerrado}

        \end{column}

      \end{columns}

      \vspace{0.35em}

      \begin{center}
        \destaca{El modelo estimará la contribución de cada medio
        publicitario manteniendo constantes los demás.}
      \end{center}

    \end{minipage}
  \end{center}

\end{frame}
% ---------------------------------------------------------------------
% Diapositiva 20 — Interpretación de los coeficientes
% ---------------------------------------------------------------------

\begin{frame}{Interpretación de los coeficientes}

  \begin{center}
    \begin{minipage}{0.90\textwidth}

      \small

      \kicker{EL APORTE ESTIMADO DE CADA VARIABLE}

      En el modelo:

      \vspace{-0.3em}

      \[
        \widehat{\text{sales}}
        =
        \beta_0
        + \beta_1(\text{TV})
        + \beta_2(\text{radio})
        + \beta_3(\text{newspaper})
      \]

      \vspace{-0.4em}

      \begin{cerrado}{Interpretación}

        \begin{itemize}
          \setlength{\itemsep}{0.3em}

          \item \(\beta_0\): ventas estimadas cuando las inversiones
                en los tres medios son iguales a cero.

          \item \(\beta_1\): cambio estimado en las ventas por cada
                unidad adicional invertida en televisión.

          \item \(\beta_2\): cambio estimado en las ventas por cada
                unidad adicional invertida en radio.

          \item \(\beta_3\): cambio estimado en las ventas por cada
                unidad adicional invertida en prensa.

        \end{itemize}

      \end{cerrado}

      \vspace{0.3em}

      \begin{center}
        \destaca{Cada cambio se interpreta manteniendo constantes
        las inversiones en los otros medios.}
      \end{center}

    \end{minipage}
  \end{center}

\end{frame}
% ---------------------------------------------------------------------
% Diapositiva 21 — Regresión simple y regresión múltiple
% ---------------------------------------------------------------------

\begin{frame}{Regresión simple y regresión múltiple}

  \begin{center}
    \begin{minipage}{0.90\textwidth}

      \small

      \kicker{DOS MODELOS, DIFERENTE CANTIDAD DE INFORMACIÓN}

      \begin{columns}[T,onlytextwidth]

        \begin{column}{0.47\textwidth}

          \begin{cerrado}{Regresión lineal simple}

            Utiliza una variable predictora:

            \[
              \hat{Y}=\beta_0+\beta_1X
            \]

            En nuestro caso:

            \[
              \widehat{\text{sales}}
              =
              \beta_0+\beta_1(\text{TV})
            \]

          \end{cerrado}

        \end{column}

        \begin{column}{0.47\textwidth}

          \begin{cerrado}{Regresión lineal múltiple}

            Utiliza dos o más variables predictoras:

            \[
              \hat{Y}
              =
              \beta_0+\sum_{j=1}^{p}\beta_jX_j
            \]

            En nuestro caso utilizaremos:

            \[
              \text{TV},\quad
              \text{radio},\quad
              \text{newspaper}
            \]

          \end{cerrado}

        \end{column}

      \end{columns}

      \vspace{0.0em}

      \begin{center}
        \destaca{Un modelo con más variables puede explicar mejor
        las ventas, pero debemos comprobarlo mediante las métricas.}
      \end{center}

      \vspace{0.0em}

      \begin{center}
        \pregunta{¿Mejorarán el \(R^2\), el MAE y el RMSE
        al incorporar las tres inversiones publicitarias?}
      \end{center}

    \end{minipage}
  \end{center}

\end{frame}
% =====================================================================
\section{Laboratorio 14}
\modolab{LABORATORIO 14}
% =====================================================================

% ---------------------------------------------------------------------
% Diapositiva 22 — Laboratorio de regresión lineal múltiple
% ---------------------------------------------------------------------

\begin{frame}{Regresión lineal múltiple con Python}

  \begin{center}
    \begin{minipage}{0.86\textwidth}

      \small

      \kicker{INCORPORAR MÁS INFORMACIÓN AL MODELO}

      En este laboratorio utilizaremos las inversiones en televisión,
      radio y prensa para estimar las ventas.

      \vspace{0.35em}

      \begin{cerrado}{Pregunta de análisis}

        \begin{center}
          \pregunta{¿Mejora la estimación de las ventas cuando
          utilizamos los tres medios publicitarios?}
        \end{center}

      \end{cerrado}

      \vspace{0.3em}

      \begin{cerrado}{Ruta de trabajo}

        \begin{enumerate}
          \setlength{\itemsep}{0.15em}

          \item Definir las tres variables predictoras.
          \item Dividir los datos en entrenamiento y prueba.
          \item Ajustar el modelo de regresión múltiple.
          \item Interpretar los coeficientes obtenidos.
          \item Evaluar y comparar los dos modelos.

        \end{enumerate}

      \end{cerrado}

    \end{minipage}
  \end{center}

\end{frame}
% ---------------------------------------------------------------------
% Cierre — Del concepto a la práctica
% ---------------------------------------------------------------------

\begin{frame}{Del concepto a la práctica}

  \begin{center}
    \begin{minipage}{0.86\textwidth}

      \kicker{LABORATORIO 13}

      En el notebook construiremos y evaluaremos diferentes modelos
      para estimar las ventas:

      \vspace{0.35em}

      \begin{cerrado}{Modelos que compararemos}

        \begin{itemize}
          \setlength{\itemsep}{0.25em}

          \item Modelo simple: \(\text{TV}\).
          \item Modelo múltiple completo:
                \(\text{TV}+\text{radio}+\text{newspaper}\).
          \item Modelo múltiple reducido:
                \(\text{TV}+\text{radio}\).

        \end{itemize}

      \end{cerrado}

      \vspace{0.35em}

      \begin{center}
        \pregunta{¿Cuál modelo ofrece el mejor equilibrio entre
        desempeño predictivo y simplicidad?}
      \end{center}

    \end{minipage}
  \end{center}

\end{frame}
\end{document}